\documentclass[12pt]{article}

\usepackage[a4paper,margin=1in]{geometry}
\usepackage{longtable}

\begin{document}

\noindent\textbf{Original Project Plan}

\begin{longtable}{|p{3cm}|p{11cm}|}
\hline
\textbf{Date} & \textbf{Planned Progress} \\
\hline

7 April &
Implementation of a minimal end-to-end experimental pipeline. This includes a simple 2D gridworld environment, generation of shortest paths using A*, training of a basic autoregressive trajectory model on the generated paths, and implementation of at least one decoding strategy to produce trajectories from the learned model. The environment will be kept deliberately simple to validate the modelling and decoding pipeline while avoiding triviality. \\
\hline

4 May &
Development of a more robust gridworld environment with configurable obstacles, start and goal states, and automated generation of shortest-path trajectories using A*. The environment should support systematic experimentation and dataset generation. \\
\hline

19 May &
Implementation of two decoding strategies (e.g., greedy decoding and stochastic sampling) and integration with the autoregressive trajectory model. Initial experiments verifying that trajectories can be generated reliably from the trained model. \\
\hline

26 May &
Preliminary evaluation of decoding strategies in the experimental environment. Initial analysis of trajectory quality and success rate. Identification of limitations of the current environment and model. \\
\hline

2 June &
Integration of the Maze2D dataset from the D4RL benchmark and adaptation of the modelling pipeline to support training and evaluation on this dataset. Verification that the autoregressive model and decoding procedures operate correctly on the benchmark data. \\
\hline

10 June &
Implementation refinements and additional experiments comparing decoding strategies. Collection of experimental results and figures to be included in the interim report. \\
\hline

16 June &
Drafting of the interim report, including description of the experimental setup, modelling approach, decoding strategies, and preliminary results. \\
\hline

20 June &
Revision and refinement of the interim report. Preparation of figures, tables, and final formatting. \\
\hline

24 June &
Final proofreading and preparation for submission of the interim report. \\
\hline

20 July &
Submission of the interim report. \\
\hline

20--24 July &
Interim presentations. \\
\hline

8 August &
Implementation of improvements to the autoregressive trajectory model to better capture trajectory structure in the dataset. Initial experiments validating the training procedure and model behaviour. \\
\hline

20 August &
Implementation and integration of additional decoding strategies, allowing systematic comparison with previously implemented decoding methods. \\
\hline

30 August &
Execution of the main experimental evaluation. This will involve systematic comparison of decoding strategies across multiple parameter settings, analysing metrics such as trajectory success rate, path optimality, diversity, and computational cost. \\
\hline

5 September &
Analysis of experimental results and preparation of figures and tables summarising the performance of the different decoding strategies. \\
\hline

10 September &
Writing and revision of the first full draft of the report describing the modelling approach, decoding strategies, experimental setup, and results obtained so far. \\
\hline

20 October &
Implementation of an application demonstrating the practical use of the proposed planning framework. This may involve applying the trained model and decoding procedures to a more complex navigation or planning scenario. Completion and refinement of the final written report will also occur during this stage. \\
\hline

30 October &
Submission of final report. \\
\hline

6 November &
Submission of the project summary document. \\
\hline

16--20 November &
Final project presentation. \\
\hline

\end{longtable}

\vspace{1em}
\noindent\textbf{Revised Project Plan}

\begin{longtable}{|p{3cm}|p{11cm}|}
\hline
\textbf{Date} & \textbf{Planned Progress} \\
\hline

7 April &
Construction of a minimal end-to-end Maze2D analysis pipeline. This includes inspection of the D4RL / Maze2D environment structure, loading and visualisation of trajectories, and initial verification of how observations, \texttt{qpos}, \texttt{qvel}, actions, and goals are represented in the benchmark data. \\
\hline

4 May &
Detailed investigation of the geometric relationship between stored Maze2D positions and maze wall layouts. Development of continuous and discrete trajectory visualisations, together with validation procedures to ensure that the plotted trajectories are consistent with the benchmark maze geometry. \\
\hline

19 May &
Construction of a discrete trajectory dataset from Maze2D-style mazes. This includes representing trajectories as alternating state and action tokens, discretising maze layouts into grid-based paths, and preparing data in a form suitable for autoregressive sequence modelling. \\
\hline

26 May &
Implementation of greedy decoding and beam search for the discrete planning setting, together with initial experiments verifying that a trained autoregressive model can generate valid trajectories from start-goal queries. \\
\hline

2 June &
Expansion from a single Maze2D example to multiple D4RL maze layouts. Construction of a benchmark dataset containing shortest valid trajectories for all connected start-goal pairs across the selected maze family. \\
\hline

10 June &
Implementation refinements to the multi-maze pipeline, including tokenisation, trajectory encoding, training-set construction, and preparation of transformer-ready inputs, labels, token-type information, and masks. \\
\hline

16 June &
Training of a small causal Transformer on the discrete trajectory data and execution of initial decoding experiments comparing greedy decoding with beam search under multiple maze layouts. \\
\hline

20 June &
Aggregation and analysis of preliminary results, including completion rate, average steps to goal, and computational-cost metrics. Preparation of figures, tables, and explanations for inclusion in the interim report. \\
\hline

24 June &
Revision and refinement of the interim report, including clarification of the probabilistic modelling formulation, dataset construction, Transformer training procedure, and decoding methodology. \\
\hline

20 July &
Submission of the interim report. \\
\hline

20--24 July &
Interim presentations. \\
\hline

8 August &
Investigation of how decoding performance changes when the trajectory model is trained on sub-optimal trajectories rather than shortest paths, testing the dependence of the current imitation-learning formulation on optimal training data. \\
\hline

20 August &
Extension of the planning framework toward an offline reinforcement learning setting, where decoding is guided by predicted reward or return rather than model probability alone. \\
\hline

30 August &
Implementation and evaluation of additional decoding strategies, such as stochastic decoding or nucleus sampling, for comparison with greedy decoding and beam search. \\
\hline

5 September &
Analysis of the expanded experimental results and preparation of additional figures and tables comparing decoding strategies under the revised settings. \\
\hline

10 September &
Writing and revision of the first full draft of the final report, consolidating the modelling formulation, experimental methodology, decoding evaluation, and discussion of results. \\
\hline

20 October &
Final refinement of the report and completion of any remaining experimental or presentation material needed for the project submission stage. \\
\hline

30 October &
Submission of final report. \\
\hline

6 November &
Submission of the project summary document. \\
\hline

16--20 November &
Final project presentation. \\
\hline

\end{longtable}

\end{document}
